\chapter{How updates work in linux}

\section{Repositories}

Lets start demystifying how linux installation works by starting with the place where our updates
come from, repositories. \textbf{Repositories} are servers that contain software that is constantly
being updated my the mantainers and where we can download our updates from. 
\\ \\
There is also the concept of \textbf{mirrors} which are literal copies of the repositories as 
they have the same content (thus the name mirrors) but are located in different places. This is
useful because you can chose mirrors that are closer to you to improve download speeds
\\ \\
If you want to see where the repository list is located in your computer, you can check 
\underbar{/etc/apt/source.list.d/}

\section{Starting the workflow: apt update}

Each distro has their own way of updating software but we'll be looking at debian based distros
for this one. The first thing we do is to ``\textbf{update}'':

\begin{lstlisting}[
    language=bash, 
    caption={Terminal},
    captionpos=t,
    label={lst:7-terminal}, 
    mathescape=true, 
    breaklines=true]
apt update
\end{lstlisting}

What ``\underbar{apt update}'' does is that it checks the list of repositories and downloads a
type of metadata known as \textbf{indexes}. These indexes contain data that apt can later use to see if
there is a pending update for a software, how to solve dependencies, where specifically do we
install the files from (inside the repository path), etc \dots
\\ \\
You see, The repositories are divided in to main paths, one known as ``\underbar{dists/}'' which contains metadata.
Using metadata is also usefull because now we don't have to check entire gigabytes of files to see
if there is an update available. Then we have the ``\underbar{pool/}'' path which contains the actual files to
be downloaded.
\\ \\
You can check the indexes at \underbar{/var/lib/apt/lists}. Usually, the files are stored in
compressed files.
\pagebreak
\section{Actually installing: apt upgrade}

Now we move into the installation workflow. ``\underbar{apt upgrade}'' starts by checking those
indexes and loading them into memory. Then it loads some sort of metadata of currently installed
packages into memory as well which can be found at \underbar{/var/lib/dpkg/status}. This means that we now can compare the metadata of the repositories
vs the ones corresponding to locally installed packages to see which require updates.
\\ \\
Something to consider is that their installation order must follow a dependency order because
we cannot install a package that depends on another one. Because of this, apt builds a dependency
graph, specifically a \textbf{Directed Acyclic Graph (DAG)}. DAGs are essentially like any other
graph that have connections with directions, and are acyclic which means that no direction can 
result in the creation of a loop.
\\ \\
Once we know which software must be updated and in which order, we can start downloading the packages.
All packages are downloaded into a cache directory at \underbar{/var/cache/apt/archives/}. This
is so if there was to be an internet outage or a package is found corrupt, the interruption wouldn't
reset the entire progress back to 0. This implies that the cache is also checked before doing
the installation.
\\ \\
Finnally, we move into the actual installation. Apt, isn't responsible for it since it cannot 
process .deb files so it calls dpkg to handle that. \textbf{dpkg} is the package installer for .deb
files. It will install all files in the cache respecting the order set by the DAG. We won't cover
how dpkg works, for now\dots

\section{The alternative: apt dist-upgrade}

Sometimes, a package will need a new dependency that wasn't contemplated. Apt by default
won't add any new packages on an update. Because of this, the packages will be labeled as ``held
back''. To force apt to install new packages and resolve all sorts of dependecy issues it my encouter,
you can use ``\underbar{apt dist-upgrade}''.

